\chapter{Drzewa, mosty i cykle fundamentalne}
\label{ch:tg-drzewa}

Drzewa są najprostszymi grafami spójnymi: można przemieszczać się
między dowolnymi dwoma wierzchołkami, ale nie ma możliwości
obchodzenia cykli. Dzięki tej sztywnej strukturze udowodnimy
wiele twierdzeń w kilku równoważnych postaciach. Później
zobaczymy, że wybór drzewa rozpinającego pozwala wyznaczać
bazę pierwszej homologii dowolnego grafu.

\section{Drzewa i ich równoważne charakteryzacje}
\label{sec:tg-def-drzewo}

\begin{definicja}[Drzewo i las]
\emph{Drzewem} nazywamy skończony, spójny, acykliczny
graf prosty. \emph{Lasem} nazywamy graf prosty
acykliczny, niekoniecznie spójny. W szczególności
każda składowa spójna lasu jest drzewem.
\end{definicja}

Dopuszczamy drzewo jednowierzchołkowe $K_1$; nie ma
ono krawędzi. Trzeba pamiętać, że acykliczność
nie oznacza braku zamkniętych spacerów (można
wracać tą samą krawędzią).

\begin{twierdzenie}[Jedyność drogi w drzewie]
Niech $G$ będzie skończonym, niepustym grafem prostym.
Wtedy $G$ jest drzewem wtedy i tylko wtedy, gdy
dla dowolnych różnych $u,v\in V(G)$ istnieje
dokładnie jedna droga prosta z $u$ do $v$.
\end{twierdzenie}

\begin{proof}
Załóżmy, że $G$ jest drzewem. Jego spójność zapewnia
istnienie spaceru między $u$ a $v$, a lemat o usuwaniu
powtórzeń --- istnienie drogi prostej. Gdyby istniały
dwie różne drogi proste $P$ i $Q$ z $u$ do $v$,
wybierzmy pierwszy wierzchołek $x$, w którym
rozchodzą się one przy przechodzeniu od $u$.
Idąc dalej po $P$, wybierzmy pierwszy wierzchołek
$y\ne x$ należący również do $Q$ po $x$;
taki wierzchołek istnieje, bo $v$ należy do obu
dróg. Fragmenty $P$ i $Q$ między $x$ i $y$
są wewnętrznie rozłączne i mają różne pierwsze
krawędzie. Ich suma tworzy cykl, co przeczy
acykliczności.

W drugą stronę, jeżeli pomiędzy każdymi dwoma
różnymi wierzchołkami jest dokładnie jedna droga
prosta, graf jest spójny. Gdyby zawierał cykl,
pomiędzy dowolnymi dwoma różnymi wierzchołkami
tego cyklu byłyby dwie drogi proste: dwa łuki
cyklu prowadzące w przeciwnych kierunkach.
Otrzymalibyśmy sprzeczność. Zatem graf jest
spójny i acykliczny, czyli jest drzewem.
\end{proof}

\begin{center}
\begin{tikzpicture}[scale=.92]
  \node[tgTealVertex] (a) at (-2.2,0) {$a$};
  \node[tgTealVertex] (b) at (-.7,1.15) {$b$};
  \node[tgTealVertex] (c) at (-.7,-1.15) {$c$};
  \node[tgTealVertex] (d) at (.8,1.5) {$d$};
  \node[tgTealVertex] (e) at (.8,.3) {$e$};
  \node[tgTealVertex] (f) at (.8,-1.35) {$f$};
  \draw[tgTree] (a)--(b)--(d) (b)--(e) (a)--(c)--(f);
  \draw[tgAccent] (d)--(b)--(a)--(c)--(f);
  \node[font=\small,text=tgNavy] at (2.3,.1) {droga $d\to f$};
\end{tikzpicture}

{\footnotesize W drzewie istnieje tylko jedna droga między
wierzchołkami $d$ i $f$ (zaznaczona na czerwono).}
\end{center}

\begin{lemat}[Liście w drzewie]
Niech $T$ będzie skończonym drzewem mającym
co najmniej dwa wierzchołki. Wtedy $T$ ma
co najmniej dwa różne wierzchołki stopnia jeden.
\end{lemat}

\begin{proof}
Wybierzmy drogę prostą $P=v_0v_1\ldots v_k$
o największej możliwej długości. Ponieważ drzewo
ma co najmniej dwa wierzchołki i jest spójne,
$k\ge1$. Pokażemy, że $d_T(v_0)=1$.
Gdyby $v_0$ miał sąsiada $w\ne v_1$, możliwe
byłyby dwa przypadki. Jeżeli $w$ leżałby na
$P$, krawędź $v_0w$ wraz z odpowiednim fragmentem
$P$ tworzyłaby cykl. Jeżeli $w$ nie leżałby
na $P$, droga $wv_0v_1\ldots v_k$ byłaby
dłuższa niż $P$. Oba przypadki są niemożliwe.
Zatem $v_0$ ma tylko jednego sąsiada $v_1$.
Tak samo $v_k$ ma stopień jeden, a ponieważ
$k\ge1$, oba końce są różne.
\end{proof}

\begin{twierdzenie}[Liczba krawędzi drzewa]
Dla każdego skończonego drzewa $T$ o $n\ge1$
wierzchołkach zachodzi $|E(T)|=n-1$.
\end{twierdzenie}

\begin{proof}
Dowodzimy przez indukcję po $n$. Dla $n=1$
drzewo nie ma krawędzi, więc wzór jest prawdziwy.
Niech $n\ge2$ i załóżmy tezę dla drzew o $n-1$
wierzchołkach. Z poprzedniego lematu $T$ ma
liść $v$ i dokładnie jedną krawędź $e$ incydentną
z $v$. Usuńmy $v$ oraz $e$. Pozostały graf $T-v$
nie ma cykli, bo jest podgrafem drzewa. Jest
też spójny: jedyna droga między dwoma dowolnymi
wierzchołkami różnymi od $v$ nie mogła przechodzić
przez liść $v$ jako wierzchołek wewnętrzny, bo
wymagałoby to dwóch różnych krawędzi przy $v$.
Stąd $T-v$ jest drzewem o $n-1$ wierzchołkach
i z założenia indukcyjnego ma $n-2$ krawędzie.
Po dołączeniu $e$ otrzymujemy $n-1$ krawędzi w $T$.
\end{proof}

\begin{wniosek}[Krawędzie lasu]
Niech $F$ będzie skończonym lasem o $n$ wierzchołkach
i $c$ składowych spójnych. Wtedy $|E(F)|=n-c$.
\end{wniosek}

\begin{proof}
Niech składowe mają odpowiednio $n_1,\ldots,n_c$
wierzchołków. Każda jest drzewem, zatem ma
$n_i-1$ krawędzi. Sumując, dostajemy
$|E(F)|=\sum_{i=1}^{c}(n_i-1)
=\sum_{i=1}^{c}n_i-c=n-c$.
\end{proof}

\begin{twierdzenie}[Charakteryzacja przez liczbę krawędzi]
Niech $G=(V,E)$ będzie skończonym grafem prostym
o $n\ge1$ wierzchołkach i $m$ krawędziach.
Następujące warunki są równoważne:
\begin{enumerate}[label=(\roman*)]
\item $G$ jest drzewem;
\item $G$ jest spójny oraz $m=n-1$;
\item $G$ jest acykliczny oraz $m=n-1$;
\item $G$ jest spójny i usunięcie dowolnej
krawędzi niszczy spójność;
\item $G$ jest acykliczny i dołączenie dowolnej
krawędzi łączącej dwa dotąd niesąsiadujące
wierzchołki tworzy cykl (jeśli takie pary istnieją).
\end{enumerate}
\end{twierdzenie}

\begin{proof}
Wykazaliśmy już (i)$\Rightarrow$(ii) oraz
(i)$\Rightarrow$(iii). Aby udowodnić
(ii)$\Rightarrow$(i), wybierzmy spośród wszystkich
spójnych podgrafów $G$ o pełnym zbiorze
wierzchołków taki, który ma najmniej krawędzi.
Nie zawiera cyklu: gdyby zawierał, usunięcie
jednej krawędzi cyklu nie zniszczyłoby spójności,
ponieważ można ją zastąpić pozostałym fragmentem
cyklu. Otrzymujemy drzewo rozpinające $T$.
Ma ono $n-1$ krawędzi; ponieważ sam $G$ także
ma $n-1$ krawędzi, musi zachodzić $E(G)=E(T)$.
Zatem $G=T$ jest drzewem.

Aby wykazać (iii)$\Rightarrow$(i), oznaczmy
przez $c$ liczbę składowych spójnych $G$.
Acykliczny $G$ jest lasem, zatem z poprzedniego
wniosku $m=n-c$. Z założenia $m=n-1$ wynika
$c=1$, a więc $G$ jest spójny i dlatego jest
drzewem.

Jeśli (i) zachodzi, to każde dwa wierzchołki
drzewa łączy jedyna droga prosta. Usunięcie
dowolnej krawędzi $uv$ niszczy tę drogę, czyli
wierzchołki $u$ i $v$ przestają być połączone;
mamy (iv). Odwrotnie, gdyby graf spełniający
(iv) zawierał cykl, usunięcie dowolnej krawędzi
tego cyklu zachowałoby spójność, co przeczy (iv).
Stąd (iv)$\Rightarrow$(i).

Na koniec załóżmy (i). Dla dwóch niesąsiadujących
wierzchołków $u,v$ istnieje w $G$ jedyna droga
prosta od $u$ do $v$. Dołączenie krawędzi $uv$
zamyka tę drogę w cykl, co daje (v).
Jeżeli natomiast (v) zachodzi, lecz $G$ nie
jest spójny, weźmy wierzchołki $u,v$ z różnych
składowych. Nie sąsiadują ze sobą, a dodanie
krawędzi $uv$ nie może utworzyć cyklu, bo
przed jej dodaniem nie było drogi między nimi.
Sprzeczność. Zatem $G$ jest spójny, a jako
acykliczny jest drzewem.
\end{proof}

\begin{uwaga}
W warunku (v) warto osobno sprawdzić przypadek grafu
pełnego. Jeżeli graf jest pełny i acykliczny,
może mieć co najwyżej dwa wierzchołki; wówczas
jest drzewem. Warunek o ,,każdej brakującej krawędzi''
pozostaje logicznie prawdziwy również wtedy, gdy
takich krawędzi nie ma.
\end{uwaga}

\section{Drzewa rozpinające}
\label{sec:tg-rozpinajace}

\begin{definicja}[Drzewo rozpinające]
Niech $G=(V,E)$ będzie grafem spójnym.
\emph{Drzewem rozpinającym} $G$ nazywamy
drugi graf $T=(V,E_T)$, gdzie $E_T\subseteq E$
oraz $T$ jest drzewem. Innymi słowy:
$T$ ma dokładnie te same wierzchołki co $G$
i pewien podzbiór jego krawędzi.
\end{definicja}

\begin{twierdzenie}[Istnienie drzewa rozpinającego]
Każdy skończony spójny graf prosty ma drzewo
rozpinające. Graf o $n$ wierzchołkach jest
spójny wtedy i tylko wtedy, gdy zawiera
dr łewo rozpinające. W szczególności każdy
spójny graf ma co najmniej $n-1$ krawędzi.
\end{twierdzenie}

\begin{proof}
Wśród wszystkich podgrafów $G$ o zbiorze
wierzchołków $V(G)$ wybieramy spójny podgraf
$T$ z minimalną liczbą krawędzi. Taki wybór
jest możliwy, bo sam $G$ jest jednym z
rozważanych podgrafów, a zbiór krawędzi jest
skończony. Gdyby $T$ zawierał cykl, można
by usunąć z niego jedną krawędź cyklu
bez utraty spójności, sprzecznie z minimalnością.
Zatem $T$ jest drzewem rozpinającym.

Jeżeli $G$ zawiera drzewo rozpinające, to
każde dwa wierzchołki łączy droga w $T$, a więc
również w $G$; jest on spójny. Ostatnia
nierówność wynika z $|E(T)|=n-1$ i z
$E(T)\subseteq E(G)$.
\end{proof}

\begin{przyklad}[Szkielet grafu z cyklami]
Dla grafu o wierzchołkach $a,b,c,d,e$ i krawędziach
$ab,bc,cd,de,ea,ac,ce$ jedno z drzew rozpinających
ma krawędzie $ab,bc,cd,de$; jest to droga
$a-b-c-d-e$. Pozostałe trzy krawędzie, po
ponownym dołączeniu, tworzą niezależne
fundamentalne cykle. Pojęcie ,,niezależności''
sformalizujemy w sekcji o przestrzeni cykli.
\end{przyklad}

\begin{center}
\begin{tikzpicture}[scale=.88]
 \node[tgBlueVertex] (a) at (-2.1,0) {$a$};
 \node[tgBlueVertex] (b) at (-.8,1.4) {$b$};
 \node[tgBlueVertex] (c) at (.9,1.4) {$c$};
 \node[tgBlueVertex] (d) at (2.1,0) {$d$};
 \node[tgBlueVertex] (e) at (0,-1.2) {$e$};
 \draw[tgGhost] (e)--(a) (a)--(c) (c)--(e);
 \draw[tgTree] (a)--(b)--(c)--(d)--(e);
 \node[font=\footnotesize,text=tgTeal] at (0,-1.78)
  {zielone: drzewo rozpinające; szare: krawędzie pozostałe};
\end{tikzpicture}
\end{center}

\begin{uwaga}[Konstrukcja algorytmiczna]
Drzewo rozpinające można otrzymać bez przeglądania
wszystkich podgrafów. Wybieramy wierzchołek
początkowy, następnie stopniowo dołączamy
krawędzie do jeszcze nieodwiedzonych wierzchołków.
Nigdy nie dołączamy krawędzi między dwiema
już odwiedzonymi pozycjami: taka krawędź
zamknęłaby cykl. Spójność gwarantuje, że
proces kończy się dopiero po odwiedzeniu
wszystkich wierzchołków. Wybór kolejności
może odpowiadać przeszukiwaniu wszerz (BFS)
lub w głąb (DFS).
\end{uwaga}

\section{Drzewo ukorzenione i odległości}
\label{sec:tg-ukorzenione}

\begin{definicja}[Korzeń, rodzic, głębokość]
\emph{Drzewem ukorzenionym} nazywamy parę
$(T,r)$, gdzie $T$ jest drzewem, a
$r\in V(T)$ jest wyróżnionym
\emph{korzeniem}. Dla $v\ne r$ jedyna droga
od $r$ do $v$ ma przedostatni wierzchołek
$p(v)$, zwany \emph{rodzicem} $v$. Każdy
wierzchołek $w$ o $p(w)=v$ jest \emph{dzieckiem}
$v$. \emph{Głębokością} $v$ nazywamy
$h(v)=d_T(r,v)$.
\end{definicja}

\begin{stwierdzenie}[Poziomy drzewa]
Niech $(T,r)$ będzie skończonym drzewem
ukorzenionym. Wtedy dla każdej krawędzi
$uv\in E(T)$ zachodzi $|h(u)-h(v)|=1$.
Każdy wierzchołek różny od $r$ ma dokładnie
jednego rodzica, a wzdłuż kolejnych rodziców
głębokość maleje o jeden aż do $r$.
\end{stwierdzenie}

\begin{proof}
Jeżeli $uv$ jest krawędzią, to
$|h(u)-h(v)|\le1$ z nierówności trójkąta
dla odległości grafowej. Gdyby głębokości
były równe, drogi od $r$ do $u$ i do $v$
musiałyby najpierw mieć wspólny odcinek,
a potem rozdzielić się (bo $u\ne v$).
Fragmenty obu dróg od ostatniego wspólnego
wierzchołka do $u,v$ oraz krawędź $uv$
utworzyłyby cykl, co jest niemożliwe.
Stąd różnica głębokości wynosi jeden.

Dla $v\ne r$ jedyna droga od $r$ do $v$
ma dokładnie jeden wierzchołek bezpośrednio
poprzedzający $v$ i jego głębokość wynosi
$h(v)-1$. Gdyby istniał inny sąsiad o
takiej głębokości, powstałyby dwie różne
drogi od $r$ do $v$, wbrew jedyności.
Kolejne zastosowanie rodzica zmniejsza
nieujemną liczbę całkowitą $h$ o jeden,
a więc kończy się dokładnie w korzeniu.
\end{proof}

\begin{center}
\begin{tikzpicture}[scale=.91]
  \node[tgRedVertex] (r) at (0,2.4) {$r$};
  \node[tgBlueVertex] (a) at (-2,1.15) {$a$};
  \node[tgBlueVertex] (b) at (1.75,1.15) {$b$};
  \node[tgTealVertex] (c) at (-3,0) {$c$};
  \node[tgTealVertex] (d) at (-1,0) {$d$};
  \node[tgTealVertex] (e) at (1,0) {$e$};
  \node[tgTealVertex] (f) at (2.55,0) {$f$};
  \draw[tgTree] (r)--(a) (r)--(b) (a)--(c) (a)--(d)
    (b)--(e) (b)--(f);
  \node[text=tgNavy,font=\footnotesize,anchor=west] at (3.25,2.4) {$h=0$};
  \node[text=tgNavy,font=\footnotesize,anchor=west] at (3.25,1.15) {$h=1$};
  \node[text=tgNavy,font=\footnotesize,anchor=west] at (3.25,0) {$h=2$};
\end{tikzpicture}

{\footnotesize Korzeń $r$ i trzy poziomy drzewa.
Każdy wierzchołek poza $r$ ma jednoznacznego rodzica.}
\end{center}

\begin{twierdzenie}[Konstrukcja BFS i odległości]
Niech $G=(V,E)$ będzie skończonym spójnym
grafem prostym, a $r\in V$. Rozważmy algorytm,
który nadaje $r$ poziom $0$, a pozostałym
wierzchołkom kolejne poziomy przez przeglądanie
sąsiadów w kolejności niemalejących poziomów.
Gdy wierzchołek $v$ zostaje pierwszy raz
odkryty jako sąsiad już odkrytego $u$,
nadajemy mu poziom $h(v)=h(u)+1$ i zapisujemy
krawędź $uv$ jako krawędź rodzicielską.
Wtedy $h(v)=d_G(r,v)$ dla każdego $v\in V$,
a krawędzie rodzicielskie tworzą drzewo
rozpinające $G$.
\end{twierdzenie}

\begin{proof}
Najpierw pokażmy, że każdemu wierzchołkowi
przypisany zostanie poziom. Gdyby po zakończeniu
procedury pozostał nieodkryty wierzchołek,
spójność grafu gwarantowałaby krawędź między
zbiorem odkrytych i nieodkrytych wierzchołków.
Algorytm musiałby ją rozważyć podczas
przeglądania wszystkich sąsiadów odkrytego
wierzchołka, więc jej drugi koniec również
zostałby odkryty. Sprzeczność.

Poziomy są rozpatrywane w kolejności $0,1,2,\ldots$.
W chwili odkrycia $v$ z poziomem $h(v)=k$
istnieje łańcuch kolejnych krawędzi
rodzicielskich długości $k$ prowadzący od $r$
do $v$. Wobec tego $d_G(r,v)\le k$.
Gdyby istniała krótsza droga długości $\ell<k$,
to indukcyjnie po kolejnych wierzchołkach tej
drogi algorytm odkryłby wierzchołek na pozycji
$j$ nie później niż przy przetwarzaniu poziomu
$j-1$. W szczególności $v$ otrzymałby poziom
co najwyżej $\ell$, sprzecznie z przypisaniem
poziomu $k$. Zatem $h(v)=d_G(r,v)$.

Każdy $v\ne r$ ma dokładnie jedną wybraną
krawędź rodzicielską. Wzdłuż tych krawędzi
poziom ściśle maleje w kierunku korzenia,
więc nie ma cykli. Możemy też dojść z
każdego wierzchołka do $r$, a więc podgraf
rodzicielski jest spójny i acykliczny.
Jest zatem drzewem rozpinającym.
\end{proof}

\begin{przyklad}[Rozbieżność między drzewem BFS a grafem]
W kwadracie $a-b-c-d-a$ wybierzmy korzeń $a$.
Mamy $d_G(a,b)=d_G(a,d)=1$ oraz $d_G(a,c)=2$.
Jedno z drzew BFS składa się z krawędzi
$ab,ad,bc$, drugie może zawierać $ab,ad,dc$.
Oba odtwarzają odległości od $a$, choć
nie są tym samym podgrafem. Nieużyta krawędź
zamyka cykl i jest celowo pomijana przez algorytm.
\end{przyklad}

\section{Mosty i przekroje krawędziowe}
\label{sec:tg-mosty}

\begin{definicja}[Most]
Niech $G$ będzie skończonym grafem prostym.
Krawędź $e\in E(G)$ nazywamy \emph{mostem},
jeżeli graf $G-e$ ma więcej składowych spójnych
niż $G$.
\end{definicja}

Usunięcie jednej krawędzi może zwiększyć liczbę składowych
co najwyżej o jeden, ponieważ jedynie składowa zawierająca
tę krawędź może ulec rozdzieleniu. W szczególności
w drzewie każda krawędź jest mostem.

\begin{twierdzenie}[Kryterium mostu]
Niech $G$ będzie skończonym grafem prostym,
a $e\in E(G)$. Krawędź $e$ jest mostem wtedy
i tylko wtedy, gdy nie należy do żadnego cyklu grafu.
\end{twierdzenie}

\begin{proof}
Niech $e=uv$ leży na cyklu $C$. Po jej
usunięciu pozostałe krawędzie cyklu wciąż
tworzą drogę z $u$ do $v$. Jeżeli spacer
w $G$ pomiędzy dowolnymi wierzchołkami
przechodził przez $e$, możemy każde takie
przejście zastąpić pozostałym fragmentem
$C$. Otrzymujemy spacer o tych samych
końcach w $G-e$. Zatem usunięcie $e$
nie zmienia składowych spójnych, więc $e$
nie jest mostem.

Załóżmy odwrotnie, że $e=uv$ nie jest
mostem. Wtedy $u$ i $v$ są połączone
w $G-e$, ponieważ w przeciwnym razie
usunięcie $e$ rozdzieliłoby ich wspólną
składową. Z lematu o redukcji spaceru
istnieje w $G-e$ droga prosta $P$ z $u$
do $v$. Ponieważ $e\notin P$, suma $P$
i krawędzi $e$ jest cyklem zawierającym $e$.
\end{proof}

\begin{definicja}[Przekrój krawędziowy]
Dla grafu $G=(V,E)$ i zbioru
$S\subseteq V$ definiujemy \emph{przekrój}
$\delta_G(S)=\{uv\in E:u\in S,\ v\in V\setminus S\}$.
Przekrój jest \emph{nietrywialny}, jeśli
$\varnothing\ne S\ne V$.
\end{definicja}

\begin{stwierdzenie}[Most jako przekrój jednoelementowy]
Niech $G$ będzie spójnym grafem prostym,
a $e\in E(G)$. Wówczas $e$ jest mostem
dokładnie wtedy, gdy istnieje niepusty
właściwy podzbiór $S\subset V(G)$ taki, że
$\delta_G(S)=\{e\}$.
\end{stwierdzenie}

\begin{proof}
Załóżmy, że $e=uv$ jest mostem. W grafie
$G-e$ wierzchołki $u,v$ należą do różnych
składowych. Niech $S$ będzie zbiorem
wierzchołków składowej zawierającej $u$.
Jest on niepusty i właściwy, a żadna krawędź
$G-e$ nie prowadzi od $S$ do jego dopełnienia.
Ponieważ $e$ łączy $u\in S$ z $v\notin S$,
mamy $\delta_G(S)=\{e\}$.

W drugą stronę, jeśli $\delta_G(S)=\{e\}$,
to po usunięciu $e$ nie istnieje żadna
krawędź łącząca $S$ z $V\setminus S$.
Oba zbiory są niepuste, więc w $G-e$
brak spaceru z dowolnego punktu jednego
zbioru do dowolnego punktu drugiego.
Graf nie jest spójny, a zatem $e$ jest mostem.
\end{proof}

\begin{center}
\begin{tikzpicture}[scale=.95]
  \node[tgBlueVertex] (a) at (-2.7,1.0) {$a$};
  \node[tgBlueVertex] (b) at (-1.1,1.0) {$b$};
  \node[tgBlueVertex] (c) at (-1.9,-.45) {$c$};
  \node[tgRedVertex] (d) at (.7,.35) {$d$};
  \node[tgRedVertex] (e) at (2.45,1.0) {$e$};
  \node[tgRedVertex] (f) at (2.45,-.55) {$f$};
  \draw[tgEdge] (a)--(b)--(c)--(a) (d)--(e)--(f)--(d);
  \draw[tgAccent] (b)--(d);
  \node[font=\small,text=tgCoral] at (-.1,1.12) {most $e$};
  \draw[tgPurple,densely dashed,line width=1pt]
     (-.18,-1.05)--(-.18,1.75);
  \node[font=\footnotesize,text=tgPurple] at (-.3,-1.35)
     {$\delta(S)=\{e\}$};
\end{tikzpicture}

{\footnotesize Jedyna krawędź między dwoma fragmentami
grafu jest mostem; krawędzie należące do trójkątów nim nie są.}
\end{center}

\begin{przyklad}[Mosty na schemacie układu połączeń]
Załóżmy, że dwa cykle $C_3$ połączono dokładnie jedną
krawędzią $e$. Krawędź $e$ jest mostem, ponieważ
nie leży na żadnym cyklu. Wszystkie sześć krawędzi
obu trójkątów leży na cyklach i nie jest mostami.
Po usunięciu $e$ graf rozpada się na dwie składowe.
\end{przyklad}

\section{Krawędzie spoza drzewa i cykle fundamentalne}
\label{sec:tg-fundamentalne}

\begin{definicja}[Cykl fundamentalny]
Niech $G$ będzie skończonym spójnym grafem
prostym i niech $T$ będzie jego drzewem
rozpinającym. Dla każdej krawędzi
$e=uv\in E(G)\setminus E(T)$ jedyna droga
$P_T(u,v)$ w $T$ wraz z krawędzią $e$
tworzy cykl $C_e$, zwany \emph{cyklem
fundamentalnym krawędzi $e$ względem $T$}.
\end{definicja}

Jeżeli graf ma $n$ wierzchołków i $m$ krawędzi,
każde drzewo rozpinające ma $n-1$ krawędzi,
a więc pozostaje dokładnie $m-n+1$ krawędzi
nie należących do drzewa. Tyle jest cykli
fundamentalnych, choć wszystkich cykli grafu
może być znacznie więcej.

\begin{center}
\begin{tikzpicture}[scale=1]
  \node[tgBlueVertex] (a) at (-2,0) {$a$};
  \node[tgBlueVertex] (b) at (-.65,1.1) {$b$};
  \node[tgBlueVertex] (c) at (.85,1.1) {$c$};
  \node[tgBlueVertex] (d) at (2.0,0) {$d$};
  \node[tgBlueVertex] (e) at (0,-1.2) {$e$};
  \draw[tgTree] (a)--(b)--(c)--(d)--(e);
  \draw[tgGhost] (a)--(e) (a)--(c);
  \draw[tgAccent] (c)--(e);
  \node[text=tgCoral,font=\footnotesize] at (1.25,-.35) {$e=ce$};
  \node[text=tgNavy,font=\footnotesize] at (0,-1.8)
    {dodanie $ce$ tworzy cykl $c-d-e-c$};
\end{tikzpicture}
\end{center}

\begin{definicja}[Przestrzeń krawędzi nad $\mathbb F_2$]
Dla skończonego grafu $G=(V,E)$ oznaczmy
przez $\mathbb F_2^E$ przestrzeń wszystkich
funkcji $x:E\to\mathbb F_2=\{0,1\}$, ze
zwykłym dodawaniem modulo $2$.
Zbiór krawędzi $F\subseteq E$ utożsamiamy
z funkcją charakterystyczną $\mathbf1_F$.
Dla dwóch zbiorów krawędzi suma odpowiada
\emph{różnicy symetrycznej}:
$F\triangle H=(F\setminus H)\cup(H\setminus F)$.
\end{definicja}

\begin{definicja}[Przestrzeń cykli]
\emph{Przestrzenią cykli} $\mathcal C(G)$ grafu
$G$ nazywamy podprzestrzeń $\mathbb F_2^E$
rozpiętą przez wektory $\mathbf1_C$,
gdzie $C$ przebiega wszystkie cykle proste grafu.
Jeśli graf nie ma cykli, przyjmujemy
$\mathcal C(G)=\{0\}$.
\end{definicja}

\begin{lemat}[Rozkład zbioru krawędzi o parzystych stopniach]
Niech $G=(V,E)$ będzie grafem prostym,
a $F\subseteq E$ takim zbiorem krawędzi,
że w podgrafie $(V,F)$ każdy wierzchołek ma
parzysty stopień. Wtedy $F$ jest rozłączną
sumą zbiorów krawędzi cykli prostych
(w szczególności $\mathbf1_F\in\mathcal C(G)$).
\end{lemat}

\begin{proof}
Dowodzimy przez indukcję po $|F|$. Gdy $F=\varnothing$,
teza jest oczywista. W przeciwnym razie
weźmy składową grafu $(V,F)$ zawierającą
co najmniej jedną krawędź. Każdy jej
nieizolowany wierzchołek ma stopień co najmniej
dwa, bo stopnie są parzyste. W skończonym
grafie o tej własności istnieje cykl:
rozpocznijmy spacer, w każdym kroku przechodząc
krawędzią inną niż ta, którą właśnie weszliśmy.
Nie możemy zakończyć ruchu w wierzchołku
stopnia co najmniej dwa; ponieważ graf jest
skończony, pewien wierzchołek zostanie
powtórzony. Fragment od pierwszego do
następnego wystąpienia powtarzającego się
wierzchołka, po ewentualnym odrzuceniu
wcześniejszego fragmentu, daje cykl prosty $C$.

Usuńmy wszystkie krawędzie $C$ z $F$.
Każdy wierzchołek na cyklu traci dokładnie
dwie krawędzie, a inne nie tracą żadnej.
Wobec tego w pozostałym podgrafie nadal
wszystkie stopnie są parzyste. Zbiór
krawędzi zmniejszył się, więc z założenia
indukcyjnego jest rozłączną sumą cykli.
Dołączając $C$, otrzymujemy rozkład całego $F$.
\end{proof}

\begin{twierdzenie}[Cykle fundamentalne tworzą bazę]
Niech $G=(V,E)$ będzie skończonym spójnym
grafem prostym, $|V|=n$, $|E|=m$,
a $T$ jego drzewem rozpinającym. Rodzina
$\{\mathbf1_{C_e}:e\in E\setminus E(T)\}$
jest bazą przestrzeni cykli $\mathcal C(G)$
nad $\mathbb F_2$. W szczególności
\[
\dim_{\mathbb F_2}\mathcal C(G)=m-n+1.
\]
\end{twierdzenie}

\begin{proof}
Najpierw pokażmy liniową niezależność.
Każdy cykl fundamentalny $C_e$ zawiera
jedną krawędź spoza drzewa, mianowicie
samą $e$; wszystkie inne jego krawędzie
należą do $T$. Dlatego w sumie
$\sum_{e\notin E(T)}a_e\mathbf1_{C_e}$
współrzędna przy krawędzi $f\notin E(T)$
jest dokładnie równa $a_f$. Jeśli cała
suma jest zerowa, każdy $a_f$ musi być
zerem. Rodzina jest więc liniowo niezależna.

Aby wykazać rozpinanie, weźmy dowolny
cykl prosty $C$ i rozważmy zbiór
$F=E(C)\triangle\bigtriangleup_{e\in E(C)\setminus E(T)} E(C_e)$.
Każda krawędź spoza drzewa należąca do $C$
pojawia się dokładnie dwa razy i znika w
różnicy symetrycznej; krawędzie spoza
drzewa nie należące do $C$ nie pojawiają się
wcale. Zatem $F\subseteq E(T)$.
Ponieważ cykl ma przy każdym wierzchołku
stopień parzysty, a własność parzystych
stopni zachowuje się przy różnicy symetrycznej,
$F$ ma wszędzie parzyste stopnie.
Ale niepusty las nie może mieć wszystkich
stopni parzystych: dowolna jego niepusta
składowa z krawędzią ma liść stopnia jeden.
Stąd $F=\varnothing$ i
$\mathbf1_C=\sum_{e\in E(C)\setminus E(T)}\mathbf1_{C_e}$.
W ten sposób każdy cykl leży w rozpiętości
cykli fundamentalnych, więc jest to baza.

Liczba jej elementów wynosi
$|E\setminus E(T)|=m-(n-1)=m-n+1$.
\end{proof}

\begin{przyklad}[Dwa cykle i ich różnica symetryczna]
Rozważmy kwadrat $a-b-c-d-a$ z przekątną $ac$.
Jako drzewo rozpinające weźmy krawędzie
$ab,bc,cd$. Krawędź $da$ tworzy fundamentalny
cykl $a-b-c-d-a$, a krawędź $ac$ tworzy
cykl $a-b-c-a$. Ich różnica symetryczna to
trójkąt $a-c-d-a$, gdyż krawędzie $ab$ i $bc$
występują w obu i znikają modulo $2$.
Mamy $n=4$, $m=5$, więc $\dim\mathcal C(G)=2$.
Wszystkie trzy cykle są zależne: suma ich
wektorów charakterystycznych wynosi zero.
\end{przyklad}

\section{Przekroje fundamentalne i parzystość przecięć}
\label{sec:tg-przekroje-fundamentalne}

\begin{definicja}[Przekrój fundamentalny]
Niech $G$ będzie skończonym spójnym grafem
prostym, a $T$ jego drzewem rozpinającym.
Dla krawędzi $f\in E(T)$ graf $T-f$
ma dokładnie dwie składowe o zbiorach
wierzchołków $S_f$ i $V\setminus S_f$.
Zbiór $D_f=\delta_G(S_f)$ nazywamy
\emph{przekrojem fundamentalnym $f$
względem $T$}.
\end{definicja}

\begin{stwierdzenie}[Parzysta liczba przecięć cyklu i przekroju]
Niech $G=(V,E)$ będzie skończonym grafem prostym,
$C$ jego cyklem prostym, a $S\subseteq V$.
Wtedy
\[
|E(C)\cap\delta_G(S)|\equiv0\pmod2.
\]
\end{stwierdzenie}

\begin{proof}
Przemieszczając się wzdłuż cyklu $C$, za
każdym razem, gdy przechodzimy krawędzią
z przekroju $\delta_G(S)$, zmieniamy
stronę podziału $S\sqcup(V\setminus S)$.
Przejście krawędzią spoza przekroju pozostawia
nas po tej samej stronie. Ponieważ spacer
jest zamknięty, po pełnym obiegu wracamy
do tej samej strony, od której zaczęliśmy.
Liczba zmian strony musi być więc parzysta.
Jest nią dokładnie liczba krawędzi
$E(C)\cap\delta_G(S)$.
\end{proof}

\begin{stwierdzenie}[Dualność fundamentalna]
Niech $G$ będzie spójnym grafem prostym,
$T$ jego drzewem rozpinającym, $e\notin E(T)$,
a $f\in E(T)$. Wtedy
\[
e\in D_f\quad\Longleftrightarrow\quad f\in C_e.
\]
Ponadto $D_f\cap E(T)=\{f\}$.
\end{stwierdzenie}

\begin{proof}
Usunięcie $f$ rozcina drzewo $T$ na
dwie składowe o zbiorach wierzchołków
$S_f$ i $V\setminus S_f$. Jedyna droga
drzewowa pomiędzy końcami $e$ przechodzi
przez $f$ dokładnie wtedy, gdy końce
$e$ leżą po przeciwnych stronach podziału.
Pierwszy warunek oznacza $f\in C_e$,
a drugi dokładnie $e\in D_f$.
W drzewie $T$ jedyną krawędzią
prowadzącą z $S_f$ do dopełnienia jest
$f$, bo gdyby była druga, po jej
dodaniu do $T-f$ powstałby cykl.
Wobec tego $D_f\cap E(T)=\{f\}$.
\end{proof}

\begin{uwaga}[Dlaczego wynik jest przygotowaniem do homologii]
W dalszej części wprowadzimy operator brzegowy
$\partial_1:\mathbb F_2^E\to\mathbb F_2^V$.
Warunek ,,każdy wierzchołek ma stopień
parzysty'' będzie równoważny $\partial_1x=0$.
Dzisiejsze twierdzenie o bazie cykli fundamentalnych
przyjmie więc postać obliczenia
$\dim\ker\partial_1=m-n+1$ dla grafu spójnego.
Po przejściu do realizacji $|G|$ otrzymamy
$H_1(|G|;\mathbb F_2)\cong\mathbb F_2^{m-n+1}$.
W szczególności $|T|$ dla drzewa $T$
nie ma nietrywialnej pierwszej homologii.
\end{uwaga}

\section{Liczba cyklomatyczna i geometryczna intuicja}
\label{sec:tg-cyklomatyczna}

\begin{definicja}[Liczba cyklomatyczna]
Niech $G$ będzie skończonym grafem prostym
o $n$ wierzchołkach, $m$ krawędziach
i $c$ składowych spójnych. Liczbę
$\beta(G)=m-n+c$ nazywamy
\emph{liczbą cyklomatyczną} grafu.
\end{definicja}

\begin{twierdzenie}
Dla każdego skończonego grafu prostego $G$
zachodzi
$\dim_{\mathbb F_2}\mathcal C(G)=m-n+c$.
W szczególności $\beta(G)\ge0$ i
$\beta(G)=0$ wtedy i tylko wtedy, gdy
$G$ jest lasem.
\end{twierdzenie}

\begin{proof}
Niech składowe spójne $G_1,\ldots,G_c$
mają odpowiednio $n_i$ wierzchołków
i $m_i$ krawędzi. Cykl nie może przebiegać
przez więcej niż jedną składową, więc
$\mathcal C(G)$ jest sumą prostą
przestrzeni $\mathcal C(G_i)$. Z twierdzenia
o bazie cykli fundamentalnych otrzymujemy
$\dim\mathcal C(G_i)=m_i-n_i+1$.
Sumując po $i$, dostajemy
$\dim\mathcal C(G)=\sum_i(m_i-n_i+1)=m-n+c$.
Wymiar przestrzeni jest nieujemny.
Jest równy zero dokładnie wtedy, gdy
żadna składowa nie zawiera cyklu, czyli
gdy cały graf jest lasem.
\end{proof}

\begin{przyklad}[Co mierzy liczba $\beta$?]
Dla drzewa na dziesięciu wierzchołkach mamy
$m=9$, $c=1$, stąd $\beta=0$.
Jeśli dołączymy do niego trzy nowe krawędzie
między dotąd niesąsiadującymi wierzchołkami,
to $n$ i $c$ pozostają bez zmian, a $m$
wzrasta o trzy. Dostajemy $\beta=3$:
nie oznacza to, że graf ma dokładnie trzy
cykle, tylko że przestrzeń cykli ma
wymiar trzy. Rzeczywistych cykli może
być znacznie więcej, ponieważ sumy
fundamentalnych cykli mogą tworzyć
nowe cykle.
\end{przyklad}

\begin{stwierdzenie}[Zmiana liczby cyklomatycznej]
Niech $G$ będzie skończonym grafem prostym,
a $e\in E(G)$. Wtedy:
\begin{enumerate}[label=(\roman*)]
\item jeżeli $e$ jest mostem, to $\beta(G-e)=\beta(G)$;
\item jeżeli $e$ nie jest mostem, to
$\beta(G-e)=\beta(G)-1$.
\end{enumerate}
\end{stwierdzenie}

\begin{proof}
Usunięcie krawędzi zmniejsza $m$ o jeden
i nie zmienia $n$. Jeśli $e$ jest mostem,
liczba składowych zwiększa się dokładnie
o jeden: $\beta(G-e)=(m-1)-n+(c+1)=m-n+c$.
Jeśli $e$ nie jest mostem, liczba
składowych pozostaje równa $c$, więc
$\beta(G-e)=(m-1)-n+c=\beta(G)-1$.
\end{proof}

\section*{Podsumowanie części podstawowej}
\addcontentsline{toc}{section}{Podsumowanie części podstawowej}

W dwóch pierwszych rozdziałach wprowadziliśmy grafy bez
odwoływania się do topologii algebraicznej. Znamy już ich
podstawowe modele, stopnie wierzchołków i macierze, wiemy,
jak budować drogi oraz rozpoznawać spójność i dwudzielność.
Rozumiemy również, dlaczego drzewa są minimalnymi grafami
spójnymi, a ich krawędzie stanowią szkielet dowolnego
spójnego grafu. Najważniejszy wynik przygotowujący dalszą
część ma postać $\dim\mathcal C(G)=m-n+c$: za chwilę
zinterpretujemy go jako obliczenie pierwszej grupy homologii.

Dalsze rozdziały rozbudują klasyczną teorię spójności,
cykle Eulera i skojarzenia. Następnie powrócimy do
macierzy incydencji, łańcuchów i operatora Laplace'a.

% ============================================================
% ITERACJA II: klasyczna teoria grafów — rozdziały 3--5
% Wszystkie rysunki w TikZ; brak osobnych zadań.
% ============================================================

